% French PDF pp. 43--48, completing the last sentence on p. 49.
\begin{proposition}{4.4.11}\label{IV.4.4.11}
Let \(R\) be a \(\widetilde{\mathcal C}\)-equivalence relation
in the sheaf \(X\). For every subsheaf \(Y\) of \(X\) stable under
\(R\), denote by \(Y'\) the quotient \(Y/R_Y\), considered as a
subsheaf of \(X'=X/R\).
Then \(Y=Y'\times_{X'}X\), and the maps \(Y\mapsto Y/R_Y\)
and \(Y'\mapsto Y'\times_{X'}X\) establish a bijective
correspondence between the set of subsheaves \(Y\) of \(X\)
stable under \(R\) and the set of subsheaves \(Y'\) of \(X'\).
\end{proposition}
\begin{proof}
If \(Y'\) is a subsheaf of \(X'\), then \(Y'\times_{X'}X\)
is a subsheaf of \(X\) stable under \(R\).
\nde{And one has \((Y'\times_{X'}X)/R=Y'\).}
If \(Y'\) is obtained by passage to the quotient from a subsheaf
\(Y\) of \(X\), then \(Y\) is a subobject of \(Y'\times_{X'}X\).
It therefore suffices to show that if one has two subsheaves
\(Y,Y_1\) of \(X\), stable under \(R\), with \(Y_1\) containing
\(Y\), and if the quotients \(Y/R_Y\), \(Y_1/R_{Y_1}\) are
identical, then \(Y=Y_1\). One is evidently reduced to proving
the same assertion in the case where \(Y_1=X\).
Then denoting by \(P\) (resp. \(Q\)) the presheaf
\(i(X)/i(R)\) (resp. \(i(Y)/i(R_Y)\)), the diagram
\[
\xymatrix{
X\ar[r] & P\\
Y\ar@{^{(}->}[u]\ar[r] & Q\ar@{^{(}->}[u] }
\]
is cartesian. Since one has a commutative diagram
\[
\xymatrix{
P\ar@{^{(}->}[r] & a(P)\\
Q\ar@{^{(}->}[u]\ar@{^{(}->}[r] & a(Q) },
\]
and since \(Q\hookrightarrow a(Q)\) is covering (4.3.11),
the monomorphism \(Q\hookrightarrow P\) is covering,
hence \(Q\) is a refinement of \(P\).
By base change, \(Y\) is a refinement of \(X\).
Since \(X,Y\) are sheaves, this implies (4.3.12) \(Y=X\).
\end{proof}

\numpara{4.4.12}\label{IV.4.4.12}
In particular, if \(Y\) is a subsheaf of \(X\), and if \(Y'=Y/R_Y\),
then the preceding correspondence defines a subsheaf \(\overline Y\)
of \(X\), stable under \(R\), containing \(Y\), and minimal
for these properties, which one calls the \emph{saturation} of
\(Y\) for the equivalence relation \(R\).

\sgasection{4.5}{The case of a topology less fine than the canonical topology}
By 4.3.6 and 4.3.8, the following conditions are equivalent
for a topology \(\mathcal T\) on \(\mathcal C\):
\begin{enumeratei}
\item \(\mathcal T\) is less fine than the canonical topology of \(\mathcal C\).
\item Every representable presheaf is a sheaf for \(\mathcal T\).
\item Every covering sieve for \(\mathcal T\) is universal effective epimorphic.
\end{enumeratei}
If \(\mathcal T\) is defined by a pretopology \(S\mapsto R(S)\),
these conditions are also equivalent to

(iv) Every family belonging to \(R(S)\) is universal effective epimorphic.

In the case where these conditions hold, the canonical functor
\(\mathcal C\to\widehat{\mathcal C}\) factors through a functor
\(j_{\mathcal C}=j:\mathcal C\to\widetilde{\mathcal C}\)
(one will also write \(j(S)=\widetilde S\)).
\nde{And one will also write \(h_S=\widehat S\), cf. the first
commutative diagram of 4.5.4.}

\begin{proposition}{4.5.1}\label{IV.4.5.1}
The functor \(j:\mathcal C\to\widetilde{\mathcal C}\) is fully
faithful and commutes with arbitrary projective limits.
In particular, it is left exact, and therefore preserves algebraic
structures defined by finite projective limits.
\end{proposition}
\begin{proof}
This follows immediately by considering the commutative diagram
\[
\xymatrix{
\mathcal C\ar[rr]\ar[dr]_j & & \widehat{\mathcal C}\\
 & \widetilde{\mathcal C}\ar[ur]_i & },
\]
and from 4.4.1 (i).
\end{proof}

Before exhibiting other properties of the functor \(j\), we must
define the topology induced on a category \(\mathcal C_{/S}\).
No longer necessarily assuming that the given topology is less
fine than the canonical topology, one does this as follows:
if \(C\) is a sieve of \(T\) in \(\mathcal C\) and if one has a
morphism \(T\to S\), then \(C\) naturally defines a sieve of
\(T\) in \(\mathcal C_{/S}\), denoted \(J(T)\)
(for the definition of a sieve of \(T\) depends only on the
category \(\mathcal C_{/T}=(\mathcal C_{/S})_{/T}\)).
If, for example, \(C\) is defined by the family \(\{T_i\to T\}\),
then its image in \(\mathcal C_{/S}\) is defined by the same
family considered as a family of morphisms of \(\mathcal C_{/S}\).
This said, the map \(T\mapsto J(T)\) defines a topology on
\(\mathcal C_{/S}\), called the topology induced by the given
topology. With the definitions of [AS], 2.3, this is the least
fine topology on \(\mathcal C_{/S}\) for which the canonical functor
\[
i_S:\mathcal C_{/S}\longrightarrow\mathcal C
\]
is a comorphism.\nde{That is, such that for every object
\(T\to S\) of \(\mathcal C_{/S}\), every covering sieve \(R'\)
of \(i_S(T\to S)=T\), considered as a sieve of
\((T\to S)\in\Ob\mathcal C_{/S}\), is covering.}
One will remark that the identifications
\[
(\mathcal C_{/S})_{/T}=\mathcal C_{/T}
\]
respect the topologies by definition.

\begin{proposition}{4.5.2}\label{IV.4.5.2}
Let \(S\) be a representable sheaf on \(\mathcal C\) and
\(F\to S\) a morphism of \(\widehat{\mathcal C}\).
For \(S'\mapsto\Hom_S(S',F)\) to be a separated presheaf
(resp. a sheaf) on \(\mathcal C_{/S}\), it is necessary and
sufficient that \(F\) be a separated presheaf (resp. a sheaf)
on \(\mathcal C\).
\end{proposition}
\begin{proof}
For every functor \(P\), one has (I 1.4.1)
\[
\Hom(P,F)=\coprod_{f\in\Hom(P,S)}\Hom_f(P,F).
\]
\nde{The original has been given in detail in what follows.}
Let \(S'\in\Ob\mathcal C\), and let \(\eta:C'\to S'\) be
a covering sieve. Since \(S\) is a sheaf, the map
\(f\mapsto f\circ\eta\) establishes a bijection
\(\Hom(S',S)\isomto\Hom(C',S)\).
Consequently the map
\[
\Hom_{\widehat{\mathcal C}}(S',F)
\longrightarrow\Hom_{\widehat{\mathcal C}}(C',F)
\]
decomposes into a \og disjoint union\fg{} of the maps
\[
\Hom_f(S',F)\longrightarrow\Hom_{f\circ\eta}(C',F).
\]
The proposition follows.
\end{proof}

\begin{corollary}{4.5.3}\label{IV.4.5.3}
The topology induced on \(\mathcal C_{/S}\) by a topology
on \(\mathcal C\) less fine than the canonical topology of
\(\mathcal C\) is less fine than the canonical topology of
\(\mathcal C_{/S}\).
\end{corollary}

\begin{corollary}{4.5.4}\label{IV.4.5.4}
Suppose that the given topology on \(\mathcal C\) is less fine
than the canonical topology. For every \(S\in\Ob\mathcal C\),
one has an equivalence of categories
\[
\widetilde{\mathcal C}_{/\widetilde S}
\longrightarrow\widetilde{\mathcal C_{/S}}.
\]
The following diagrams are commutative up to isomorphism
(all unlabeled arrows are equivalences):
\[
\xymatrix@C=1.4em@R=3em{
\mathcal C_{/S}\ar[r]^{(j_{\mathcal C})_{/S}}\ar@{=}[d]
 & \widetilde{\mathcal C}_{/\widetilde S}
\ar[r]^{(i_{\mathcal C})_{/\widetilde S}}\ar[d]_{\simeq}
 & \widehat{\mathcal C}_{/\widehat S}
\ar[r]^{(a_{\mathcal C})_{/\widehat S}}\ar[d]_{\simeq}
 & \widetilde{\mathcal C}_{/\widetilde S}\ar[d]_{\simeq}\\
\mathcal C_{/S}\ar[r]_{j_{\mathcal C_{/S}}}
 & \widetilde{\mathcal C_{/S}}\ar[r]_{i_{\mathcal C_{/S}}}
 & \widehat{\mathcal C_{/S}}\ar[r]_{a_{\mathcal C_{/S}}}
 & \widetilde{\mathcal C_{/S}} },
\]
\[
\xymatrix@C=1.3em@R=3em{
(\widetilde{\mathcal C}_{/\widetilde S})_{/\widetilde T}
\ar[r]\ar[d]_{\simeq}
 & (\widetilde{\mathcal C_{/S}})_{/\widetilde T}
\ar[dr] & \\
\widetilde{\mathcal C}_{/\widetilde T}
\ar[r] & \widetilde{\mathcal C_{/T}}\ar[r]_{\simeq}
 & \widetilde{(\mathcal C_{/S})_{/T}} }.
\]
\end{corollary}
\begin{proof}
The commutativity of the first two squares follows from the
definition of the equivalence
\(\widetilde{\mathcal C}_{/\widetilde S}
\to\widetilde{\mathcal C_{/S}}\).
To prove the commutativity of the last, one must see that
the following square is commutative:
\[
\xymatrix{
\widehat{\mathcal C}_{/\widehat S}\ar[r]^{L'}\ar[d]
 & \widehat{\mathcal C}_{/\widehat S}\ar[d]\\
\widehat{\mathcal C_{/S}}\ar[r]_{L''}
 & \widehat{\mathcal C_{/S}} },
\]
where \(L'\) is the restriction of the functor \(L_{\mathcal C}\)
to \(\widehat{\mathcal C}_{/\widehat S}\), and \(L''\) the
functor \(L_{\mathcal C_{/S}}\); this is easily seen by returning
to the definition of the functors \(L\) (given after 4.3.9).
As for the second diagram, it is simply the restriction to
the categories of sheaves of the corresponding diagram for
the categories of presheaves (Exposé I, no.\ 1), which is commutative.
\end{proof}

\begin{scholium}{4.5.5}\label{IV.4.5.5}
The various assertions of this number show that in the case
where the given topology is less fine than the canonical topology,
one may identify \(\mathcal C\) with a full subcategory of
\(\widetilde{\mathcal C}\), itself a full subcategory of
\(\widehat{\mathcal C}\), and that in this identification
one may indulge in the customary abuses of language concerning
\(\mathcal C\to\widehat{\mathcal C}\), justified by the preceding
commutativities. Let us remark explicitly that the first diagram
of 4.5.4 shows that one may use the functor \(a\) without special precautions.
\end{scholium}

We shall see in the next number that the identification of
\(\mathcal C\) with a full subcategory of \(\widetilde{\mathcal C}\)
(contrary to what happened for \(\widehat{\mathcal C}\)) commutes
with the formation of certain inductive limits, and we shall
then say how to use this fact.

From now on, unless expressly stated otherwise, we shall assume
that the given topology is less fine than the canonical topology,
and we shall systematically make the identifications set out above.

\begin{proposition}{4.5.6}\label{IV.4.5.6}
Let \(F,G\) be two sheaves over \(S\), and \(f:F\to G\) an
\(S\)-morphism. The following conditions on \(f\) are equivalent:

(i) \(f\) is a monomorphism (resp. an epimorphism, resp. an
isomorphism) in \(\widetilde{\mathcal C}\).

(ii) \(f\) is a monomorphism (resp. an epimorphism, resp. an
isomorphism) in
\(\widetilde{\mathcal C}_{/S}=\widetilde{\mathcal C_{/S}}\).
\end{proposition}
\begin{proof}
For monomorphisms and isomorphisms this is evident (it is a
question of presheaves). For epimorphisms, this follows from
the description of epimorphisms as covering morphisms (4.4.3)
and the fact that, by definition of the induced topology,
these are the same in \(\widetilde{\mathcal C}\) and
\(\widetilde{\mathcal C}_{/S}\).
\end{proof}

\begin{proposition}{4.5.7}\label{IV.4.5.7}
Let \(f:F\to G\) be a morphism of sheaves.
The following conditions are equivalent:

(i) \(f\) is a monomorphism (resp. an epimorphism, resp. an isomorphism).

(ii) For each \(S\in\Ob\mathcal C\), \(f_S:F_S\to G_S\)
is locally a monomorphism (resp. an epimorphism, resp. an
isomorphism), that is:

(iii) For each \(S\in\Ob\mathcal C\), the set of \(T\to S\)
such that \(F_T\to G_T\) is a monomorphism (resp. an
epimorphism, resp. an isomorphism) is a refinement of \(S\).

If the given topology is defined by a pretopology \(R\),
these conditions are also equivalent to the following:

(iv) For each \(S\in\Ob\mathcal C\), there exists a covering
family \(\{S_i\to S\}\in R(S)\) such that for every \(i\),
\(F_{S_i}\to G_{S_i}\) is a monomorphism (resp. an
epimorphism, resp. an isomorphism).

If the category \(\mathcal C\) possesses a final object \(e\),
one may restrict oneself to taking \(S=e\) in conditions
(ii), (iii), (iv).
\end{proposition}
\begin{proof}
One evidently has (ii) \(\Leftrightarrow\) (iii)
\(\Leftrightarrow\) (iv).
To prove the equivalence of (i) and (ii), as well as the
supplement concerning the final object, one must show that
(ii) implies (i) and that the notions considered are stable
under base extension. Let us first prove the latter point.
For monomorphisms and isomorphisms, this is evident
(it is a question of presheaves). For epimorphisms, it follows
from the fact that every epimorphism of sheaves is universal (4.4.3).

Finally let us show that (ii) implies (i).
First suppose that \(f_S:F_S\to G_S\) is locally a monomorphism
(resp. an isomorphism). There then exists a covering sieve \(C\)
of \(S\) such that for every \(T\to C\), \(f_T\) is a
monomorphism (resp. an isomorphism). Since a projective limit
of monomorphisms (resp. isomorphisms) is one,
\(\Hom(C,f)\) will be a monomorphism (resp. an isomorphism)
(cf. 4.1.4). The commutative diagram
\[
\xymatrix{
\Hom(C,F)\ar[r]^{\Hom(C,f)} & \Hom(C,G)\\
F(S)\ar[u]^{\simeq}\ar[r]_{f(S)} & G(S)\ar[u]_{\simeq} }
\]
then shows that \(f(S)\) is injective (resp. bijective).
Finally suppose that \(f:F\to G\) is locally an epimorphism,
and let \(H\subset G\) be its image.
For each \(S\to G\), \(H\times_G S\) is the image of
\(f\times_G S:F\times_G S\to S\).
To show that \(f\) is an epimorphism, one must show that
\(H\) is a refinement of \(G\), that is, \(H\times_G S\)
is a refinement of \(S\) for each \(S\).
But since this is so after every base change \(T\to S\)
from a refinement of \(S\) (since \(f\times_G S\) is locally
covering), \(H\times_G S\) is indeed a refinement of \(S\)
(Axiom (T 2)).
\end{proof}

\begin{corollary}{4.5.8}\label{IV.4.5.8}
Let \(F,G\) be two sheaves over \(S\), and \(f:F\to G\)
an \(S\)-morphism. For \(f\) to be a monomorphism, resp.
an epimorphism, resp. an isomorphism, it is necessary and
sufficient that it be so locally on \(S\).
\end{corollary}
\begin{remark}{4.5.9}\label{IV.4.5.9}
The proof of the proposition shows that it remains valid,
for the part concerning monomorphisms (resp. isomorphisms),
when one assumes only that \(F\) is a separated presheaf
(resp. a sheaf) and \(G\) is any presheaf
(resp. a separated presheaf).
\end{remark}
Let us temporarily return to the case of an arbitrary topology
and give a definition.
\begin{definition}{4.5.10}\label{IV.4.5.10}
Let \(G\to F\) be a morphism of \(\widehat{\mathcal C}\).
One says that \(G\) is a \emph{relative sheaf over \(F\)} if
for each \(F\)-functor \(H\) and each refinement \(R\) of \(H\),
the canonical map
\[
\Hom_F(H,G)\isomto\Hom_F(R,G)\tag{+}
\]
is bijective.
\end{definition}
Proposition 4.5.2 generalizes at once:
\begin{proposition}{4.5.11}\label{IV.4.5.11}
If \(F\) is a sheaf, \(G\) is a relative sheaf over \(F\)
if and only if it is a sheaf.
\end{proposition}
\begin{lemma}{4.5.12}\label{IV.4.5.12}
In the situation \(X\to T\to S\) (where \(X,T,S\) are three
objects of \(\widehat{\mathcal C}\)), if \(X\) is a relative
sheaf over \(T\), then \(U=\prod_{T/S}X\) is a relative sheaf over \(S\).
\end{lemma}
\begin{proof}
Indeed, for every \(S\)-functor \(Y\), one has
\[
\Hom_S(Y,U)=\Hom_T(T\times_S Y,X).
\]
If \(C\) is a sieve of \(Y\), then \(T\times_S C\) is a sieve
of \(T\times_S Y\); one concludes at once.
\end{proof}
\begin{corollary}{4.5.13}\label{IV.4.5.13}
The presheaves \(\uHom_{T/S}(X,Y)\),
\(\underline{\Isom}_S(X,Y)\), etc., are sheaves when the
arguments which occur in them are also sheaves.
\end{corollary}
\begin{proof}
Indeed, all these presheaves are constructed with the aid of
fibered products and presheaves \(\prod\) (I 1.7 and II 1).
It therefore suffices to verify the result for a presheaf
\(\prod_{T/S}X\); in this case, the assertion follows from
4.5.11 and 4.5.12.
\end{proof}

\sgasection{4.6}{Description of the quotient of a sheaf by an equivalence relation}
Recall that we assume that the given topology \(\mathcal T\)
is less fine than the canonical topology.
\begin{proposition}{4.6.1}\label{IV.4.6.1}
Let \(R\xrightarrow{p_1,p_2}X\) be a
\(\widetilde{\mathcal C}\)-equivalence relation in the sheaf \(X\).
Let \(F\in\Ob\widehat{\mathcal C}\) be defined as follows:
for each \(S\) of \(\mathcal C\),
\[
F(S)=\left\{\begin{array}{l}
\text{sub-}S\text{-sheaves }Z\text{ of }X_S\text{ stable under }R\times S,\\
\text{whose quotient by }R_Z\text{ is }S,\text{ that is, such that}\\
\text{the diagram }R_Z\rightrightarrows Z\to S\text{ is exact}
\end{array}\right\}.
\]
\nde{\(R\times S\) denotes the equivalence relation in \(X\times S\)
defined by \(R\times S_{\mathrm{diagonal}}\subset X\times X\times S\times S\),
and \(R_Z\) is the equivalence relation which it induces in \(Z\)
(cf. 3.1.6).}
Then, for every sheaf \(Y\), \(\Hom(Y,F)\) is identified with the set:
\[
\left\{\begin{array}{l}
\text{sub-}Y\text{-sheaves of }X\times Y\text{ stable under }R\times Y\\
\text{and whose quotient is }Y
\end{array}\right\}.
\]
In particular, the subsheaf \(R\) of \(X\times X\) corresponds
to an element \(p\) of \(\Hom(X,F)\), and the diagram
\[
\xymatrix{R\ar@<.5ex>[r]^{p_1}\ar@<-.5ex>[r]_{p_2}
 & X\ar[r]^p & F }
\]
is exact, hence identifies \(F\) with the quotient sheaf \(X/R\).
\end{proposition}
\begin{proof}
Indeed, set \(Q=X/R\). For every sheaf \(Y\) and every morphism
\(f\in\Hom(Y,Q)\) corresponding to a section
\(s:Y\hookrightarrow Q\times Y\), consider the diagram
\[
\xymatrix{
R\times Y\ar@<.5ex>[r]\ar@<-.5ex>[r] & X\times Y\ar[r] & Q\times Y\\
 & Z\ar@{^{(}->}[u]\ar[r] & Y\ar@{^{(}->}[u]_s }.\tag{*}
\]
Here the square is cartesian. It is immediate by 4.4.11 that
\(Z\) is a sub-\(Y\)-sheaf of \(X\times Y\), stable under
\(R\times Y\), whose quotient is \(Y\), and that conversely,
every \(Z\) of this type comes from a unique section of
\(Q\times Y\) over \(Y\).
% Proof continues in en-09.tex.
